% ===================== PREAMBULE COMMUN (à copier VERBATIM en tête de CHAQUE .tex) =====================
\documentclass[11pt,a4paper]{article}
\usepackage[utf8]{inputenc}
\usepackage[T1]{fontenc}
\usepackage{lmodern}
\usepackage[french]{babel}
\usepackage{amsmath,amssymb,mathtools}
\usepackage{icomma}
\usepackage{xcolor}
\usepackage{colortbl}
\usepackage{tikz}
\usepackage{pgfplots}
\usepackage[most]{tcolorbox}
\usepackage{enumitem}
\usepackage{array}
\usepackage{booktabs}
\usepackage{multicol}
\usepackage{fancyhdr}
\usepackage{caption}
\captionsetup{labelformat=empty,justification=centering,font=small}
\usepackage[hidelinks]{hyperref}
\usetikzlibrary{arrows.meta,positioning,calc,angles,quotes,patterns,backgrounds,shapes.geometric}
\pgfplotsset{compat=1.18}
\usepackage[a4paper,margin=2.2cm,headheight=26pt,headsep=10pt]{geometry}

\definecolor{primary}{RGB}{219,54,77}
\definecolor{primarydark}{RGB}{150,28,46}
\definecolor{primarylight}{RGB}{250,224,228}
\definecolor{defcol}{RGB}{0,114,178}
\definecolor{thmcol}{RGB}{0,140,120}
\definecolor{excol}{RGB}{0,135,90}
\definecolor{methcol}{RGB}{90,90,100}
\definecolor{courbe}{RGB}{40,90,200}
\definecolor{courbeB}{RGB}{210,60,40}
\definecolor{courbeC}{RGB}{30,150,80}
\definecolor{axiscol}{RGB}{60,60,60}
\definecolor{gridcol}{RGB}{200,205,215}

\tcbset{boxbase/.style={enhanced,breakable,boxrule=0.9pt,arc=2.5pt,
   left=3mm,right=3mm,top=2.3mm,bottom=2.3mm,
   fonttitle=\bfseries,coltitle=white,toptitle=0.6mm,bottomtitle=0.6mm}}
\newtcolorbox{methode}[1][]{boxbase,colback=methcol!7,colframe=methcol,sharp corners,
  title={\textsc{Méthode}\ifx\relax#1\relax\relax\else\textmd{ --- \itshape #1}\fi}}
\newtcolorbox{exemple}[1][]{boxbase,colback=excol!5,colframe=excol,
  title={\textsc{Exemple}\ifx\relax#1\relax\relax\else\textmd{ : \itshape #1}\fi}}
\newtcolorbox{idee}[1][]{boxbase,colback=defcol!6,colframe=defcol,
  title={\textsc{L'essentiel}\ifx\relax#1\relax\relax\else\textmd{ : \itshape #1}\fi}}
\newtcolorbox{piege}[1][]{boxbase,colback=primary!7,colframe=primary,
  title={\textsc{Piège}\ifx\relax#1\relax\relax\else\textmd{ : \itshape #1}\fi}}

\pgfplotsset{repere/.style={axis lines=middle,
    axis line style={-{Stealth[length=2.2mm]},axiscol,line width=0.8pt},
    label style={font=\footnotesize},tick label style={font=\scriptsize},
    every axis x label/.style={at={(current axis.right of origin)},anchor=north west},
    every axis y label/.style={at={(current axis.above origin)},anchor=south east},
    xlabel={$x$},ylabel={$y$},grid=both,minor tick num=0,
    major grid style={gridcol,line width=0.4pt},samples=200,clip=false}}

\pagestyle{fancy}\fancyhf{}
\fancyhead[L]{\small\textcolor{primarydark}{\textbf{Fiche 4} — Les fonctions}}
\fancyhead[R]{\small\textcolor{primarydark}{\textsc{Carnet d'entraînement}}}
\fancyfoot[C]{\small\thepage}
\renewcommand{\headrulewidth}{0.8pt}
\renewcommand{\headrule}{\hbox to\headwidth{\color{primary}\leaders\hrule height \headrulewidth\hfill}}
\usepackage{titlesec}
\titleformat{\section}{\Large\bfseries\color{primarydark}}{\thesection}{0.6em}{}
\titleformat{\subsection}{\large\bfseries\color{primary}}{\thesubsection}{0.6em}{}

\newcommand{\R}{\mathbb{R}}\newcommand{\Z}{\mathbb{Z}}\newcommand{\N}{\mathbb{N}}
\newcommand{\Df}{D_f}
\newcommand{\croit}{\textcolor{thmcol}{\Large$\nearrow$}}
\newcommand{\decroit}{\textcolor{thmcol}{\Large$\searrow$}}

\newcommand{\exercice}[1]{\medskip\noindent{\color{primarydark}\bfseries\large Exercice #1}\par\nopagebreak\smallskip}
\newcommand{\titrecarnet}[3]{%
\begin{center}\begin{tikzpicture}
\node[rectangle,rounded corners=7pt,fill=primary,text=white,minimum width=15cm,
  minimum height=1.7cm,align=center,inner sep=8pt]
  {{\LARGE\bfseries #1}\\[2pt]{\normalsize #2}};
\end{tikzpicture}\end{center}
\begin{center}\itshape\color{primarydark}#3\end{center}\medskip}

% ---- RÈGLES D'USAGE DES MACROS ----
% * \croit et \decroit s'emploient en mode TEXTE (dans une cellule de tableau),
%   JAMAIS entre $...$ (ils contiennent déjà un $\nearrow$).
% * \titrecarnet{Titre}{Sous-titre}{Ligne en italique} : bandeau de titre.
% * \exercice{1} : titre d'exercice.
% * Boîtes : \begin{idee}...\end{idee}, \begin{exemple}[titre]...\end{exemple},
%   \begin{methode}[titre]...\end{methode}, \begin{piege}[titre]...\end{piege}.
% Le corps du document commence après \begin{document}.
\begin{document}
\titrecarnet{Thème 5 — Trigonométrie, exponentielle \& logarithme}{Corrigé — niveau Difficile}{Détail des calculs.}

\exercice{1}
Doubler tous les $20$ jours signifie multiplier par $2$ à chaque palier de $20$ jours.
\begin{enumerate}[label=\alph*),leftmargin=1.8em,topsep=2pt]
  \item $8=2^3$ : il faut \textbf{3 doublements}, soit $3\times 20=\boxed{60\ \text{jours}}$.
  \item $16=2^4$ : \textbf{4 doublements}, soit $4\times 20=\boxed{80\ \text{jours}}$.
  \item $\boxed{N(t)=N_0\cdot 2^{\,t/20}}$ ($t$ en jours). Vérification : $N(20)=2N_0$, $N(60)=2^3 N_0=8N_0$. \checkmark
\end{enumerate}

\exercice{2}
\begin{enumerate}[label=\alph*),leftmargin=1.8em,topsep=2pt]
  \item Avec $\cos^2 x+\sin^2 x=1$ et $\cos x>0$ (premier quadrant) :
        \[ \cos x=\sqrt{1-\sin^2 x}=\sqrt{1-\Big(\tfrac34\Big)^2}=\sqrt{1-\tfrac{9}{16}}=\sqrt{\tfrac{7}{16}}=\boxed{\dfrac{\sqrt7}{4}}. \]
  \item $\sin(2x)=2\sin x\cos x=2\times\dfrac34\times\dfrac{\sqrt7}{4}=\boxed{\dfrac{3\sqrt7}{8}}$.
\end{enumerate}

\exercice{3}
\begin{enumerate}[label=\alph*),leftmargin=1.8em,topsep=2pt]
  \item Dans le troisième quadrant $\cos\alpha<0$. Avec $\cos^2\alpha=1-\sin^2\alpha=1-(-0{,}8)^2=1-0{,}64=0{,}36$ :
        \[ \cos\alpha=-\sqrt{0{,}36}=\boxed{-0{,}6}. \]
  \item $\tan\alpha=\dfrac{\sin\alpha}{\cos\alpha}=\dfrac{-0{,}8}{-0{,}6}=\boxed{\dfrac43}$.
\end{enumerate}

\exercice{4}
\textbf{Condition d'existence :} $x>0$ \emph{et} $x-2>0$, donc $x>2$.
En regroupant : $\log_2\big(x(x-2)\big)=3\Leftrightarrow x(x-2)=2^3=8$, soit
\[ x^2-2x-8=0\Leftrightarrow (x-4)(x+2)=0\Leftrightarrow x=4 \text{ ou } x=-2. \]
Seul $x=4$ vérifie $x>2$. Donc $\boxed{x=4}$ (la solution $x=-2$ est rejetée).

\exercice{5}
\begin{enumerate}[label=\alph*),leftmargin=1.8em,topsep=2pt]
  \item Après $12$~h : $\dfrac{Q(12)}{Q_0}=\Big(\dfrac12\Big)^{12/6}=\Big(\dfrac12\Big)^2=\boxed{\dfrac14}$ (deux demi-vies).
  \item On cherche $t$ tel que $\Big(\dfrac12\Big)^{t/6}=\dfrac18=\Big(\dfrac12\Big)^3$, d'où $\dfrac{t}{6}=3\Rightarrow t=\boxed{18\ \text{heures}}$ (trois demi-vies).
\end{enumerate}
\end{document}
